\section{Enhanced NRB \& ORB (ENRB \& EORB)}
\label{sec:enh}

\subsection{Evaluation criteria}
\label{sec:enh_criteria}

The evaluation criteria for \acs{enrb} and \acs{eorb} are the same as for the \acs{nrb} and \acs{orb} products, since the underlying radiometric, geometric, and statistical properties are preserved. The additional sidelobe suppression or ghost removal improves visual quality and noise characteristics but does not alter the fundamental assessment dimensions, such as radiometric consistency, slope-effect removal, \acs{dem} and orbital dependencies, or the statistical stability of $\sigma_0$ values. Therefore, \acs{enrb}/\acs{eorb} inherit the same evaluation framework and quality benchmarks as their parent \acs{nrb}/\acs{orb} products.

\subsection{Theoretical models}
\label{sec:enh_models}

The \acs{enrb}/\acs{eorb} products' theoretical models to evaluate performances are directly inherited from the \acs{nrb} and \acs{orb} products, as their generation is based on the same processing principles. Any performance assessment therefore follows the same criteria and methodologies.

\subsection{Computational performance}
\label{sec:enh_computation}

The computational performance of \acs{enrb} and \acs{eorb} directly follows that of the \acs{nrb} and \acs{orb} products. Since their generation relies on the same processing chain, including rasterization, shadow masking, and possible stack-based co-registration, the overall cost is essentially inherited. The only additional burden comes from sidelobe suppression or ghost removal, which is marginal compared to the dominant steps already present in \acs{nrb}/\acs{orb} processing. Therefore, \acs{enrb}/\acs{eorb} maintain the same scalability and efficiency characteristics as their parent products.

\subsection{Criticalities and weakness}
\label{sec:enh_criticalities}
For \acs{enrb} products, the same criticalities and weaknesses described for \acs{nrb} apply. Since \acs{enrb} products are derived directly from \acs{nrb}, their performance is equally dependent on the accuracy of orbital data, \acs{dem} quality, and timing corrections. The main difference lies in the application of filtering techniques (ghost removal and sidelobe suppression), which may slightly mitigate noise but do not alter the fundamental dependencies of the product.

For \acs{eorb} products, the situation mirrors that of \acs{orb}. The reliability of the output is still determined by the quality of the input \acs{enrb} data, wind fields, and the geophysical model function, while the filtering step has only a marginal effect. Therefore, the same considerations regarding reference wind choice and angle of incidence compensation remain valid.
